% ============================================================
% §6 Conclusion — OR journal style: full section
% Summary, managerial insights, limitations, future directions
% ============================================================
\section{Conclusion}
\label{sec:conclusion}

% --- Summary ---
This paper addresses inference-time search computation allocation across
the denoising trajectory of diffusion models under a strict budget of function
evaluations.
By separating local noise refinement from global timestep scheduling, we
introduce a unified two-level framework that subsumes existing inference-time
search methods as special cases.
Within this framework, we develop GAINS, a global scheduling algorithm that
combines offline sensitivity profiling with online feedback control.
Offline profiling learns a coarse per-step budget from calibration
runs, encoding model-dependent patterns of timestep importance.
Online control refines this budget on a per-instance basis using windowed
gain and variance signals for early stopping, while enforcing an exact NFE
constraint.

% --- Key findings ---
Our theoretical analysis justifies both design components.
Under a location-scale model of per-step score distributions, optimal offline
allocation follows a water-filling structure that assigns more
compute to high-variance timesteps (\cref{prop:offline}), and any fixed
allocation incurs a Jensen gap that online adaptation can partially recover
(\cref{prop:online}).
Experiments on Stable Diffusion, EDM, and flow-based models confirm
these predictions: GAINS reaches matched quality with 20--50\% fewer
function evaluations, and the gains persist
across different local search operators, prompt sets, and model architectures.

% --- Managerial insights / operational protocols ---
\paragraph{Operational guidelines.}
Three concrete protocols follow from our results.
(i)~\emph{Always allocate non-uniformly.}  Non-uniform compute distribution across
denoising timesteps strictly dominates uniform distribution whenever timestep
sensitivities differ, a condition satisfied by every model we tested,
improving quality by 3--4\% at matched NFE or reducing NFE by 20--50\% at matched quality;
practitioners should run the one-time profiling step before deploying any diffusion-based generation pipeline.
(ii)~\emph{Amortize profiling cost.}  The offline profiling overhead is modest (50--100 calibration samples, amounting to roughly 1--2\% of the compute spent on a typical generation campaign) and is paid once per model checkpoint;
all subsequent inference calls reuse the same budget schedule.
(iii)~\emph{Swap local search independently.}  The two-level design is modular:
GAINS pairs with any local search operator, so teams can upgrade the
local method (e.g., from random search to Langevin refinement) without
retuning the global schedule.

% --- Limitations ---
\paragraph{Limitations.}
Three limitations bound the current scope.
First, our theoretical analysis assumes location-scale per-step score
distributions; heavier tails or multimodal score distributions would require
extended analysis.
Second, the offline profiling stage treats timestep sensitivities as
approximately stationary across prompts; highly heterogeneous prompt
distributions may need finer-grained, prompt-conditional profiling.
Third, our experiments use two lightweight verifiers (Brightness and
Compressibility); testing GAINS with richer verifiers (e.g.,
human-preference models, domain-specific quality metrics) would broaden the
empirical scope.

% --- Future directions (integrated into limitations, per Rule 39) ---
Addressing these limitations opens natural extensions.
The MDP formulation in \cref{sec:mdp} admits richer policy classes,
including backtracking, operator mixing across timesteps, and
multi-fidelity verifiers.
These extensions may further push the quality-compute frontier but
likely require learned optimization.
Beyond diffusion models, the core principle of non-uniform compute
allocation guided by sensitivity profiling applies wherever computation
is a scarce resource in sequential decisions.
Two immediate candidates are adaptive simulation budgeting in stochastic
optimization and dynamic evaluation-effort distribution across stages of
multi-stage decision problems~\citep{chen2000simulation}.
